\documentclass[11pt,a4paper]{article}
\usepackage[T1]{fontenc}
\usepackage[utf8]{inputenc}
\usepackage[a4paper,top=0.6in,bottom=0.6in,left=0.8in,right=0.8in]{geometry}
\usepackage{mathptmx}
\usepackage{setspace}
\usepackage{enumitem}
\usepackage{array}
\usepackage{titlesec}
\usepackage{parskip}

\setstretch{1.0}
\setlength{\parskip}{3pt}
\setlist[itemize]{leftmargin=1.2em,itemsep=0.05em,topsep=0.15em,parsep=0pt}
\titleformat{\section}{\normalfont\normalsize\bfseries}{\thesection}{0.7em}{}
\titlespacing*{\section}{0pt}{0.9ex}{0.45ex}
\pagestyle{empty}

\begin{document}

\begin{center}
{\large\bfseries How A.R.T.E.M.I.S Differs From the Closest Published Work}\\[0.3em]
{\small F26-008 \textbar{} Hashir Rauf (23L-2572), Hira Khalid (23L-2594),
Syeda Doureesha Batool (23L-2651) \textbar{} Supervisor: Mr.\ Mubasher Hussain}
\end{center}
\vspace{-0.4em}
\hrule
\vspace{0.6em}

\section{The Two Closest Papers}

\textbf{HULA} (Takerngsaksiri et al., ICSE-SEIP 2025, IEEE) deploys a
plan-then-approve-then-code agent at Atlassian. Engineers approved the agent's
plan in 82 percent of cases, establishing that human approval of an explicit
plan is a workable interaction model rather than an obstacle.

\textbf{SoK: Attack and Defense Landscape of Agentic AI Systems} (Kim, Guo and
Song, USENIX Security 2026) systematises 51 attacks and 60 defences across 128
papers. It concludes that architectural confinement is more dependable than
detecting malicious input, and documents that human approval degrades through
decision fatigue.

Together these two represent the state of the art on the two axes our project
occupies: the approval loop as an interaction model, and architectural
confinement as a security position.

\section{What We Take From Them}

Our system adopts the plan-and-approve loop that HULA validates, and the
confinement-over-detection position that the SoK recommends. Neither is claimed
as our contribution.

\section{Where We Differ}

\begin{itemize}
  \item \textbf{Scope is granted by the user, not assumed.} HULA's agent
        operates across a repository and the surveyed systems generally assume
        broad machine access. Our system resolves every path through one
        component that canonicalises before testing containment, so a file
        outside a folder the user explicitly granted is unreachable, not merely
        refused. The SoK argues for this confinement; we apply it to a
        user-chosen folder on an ordinary desktop.
  \item \textbf{Destruction is unavailable rather than guarded.} There is no
        approval dialogue for deleting a file because there is no code path to
        it. A gate the user can click through is a policy; a gate that does not
        exist is an architecture.
  \item \textbf{Reversal exists.} Neither paper provides a facility to undo an
        action once executed. We record how to reverse each change before
        making it, and reverse an approved plan as the single unit the user
        approved rather than as separate operations.
  \item \textbf{The user is not assumed to be an engineer.} HULA's reviewers
        are professional engineers; the SoK notes that most defences target
        organisations with security staff. Our previews are written in plain
        language for a student or ordinary knowledge worker, and our research
        with 374 respondents is drawn from that population.
  \item \textbf{Approval is graduated, not uniform.} Actions are classified into
        four risk tiers, so reads never prompt, reversible changes prompt
        according to user configuration, and actions that leave the machine
        always prompt. This directly targets the fatigue the SoK documents.
  \item \textbf{The system discloses what it holds.} A panel reports every
        file, fact and action retained, and every entry can be deleted. The SoK
        notes that agent memory is typically an opaque process the user cannot
        inspect.
  \item \textbf{Beyond code.} HULA addresses software tasks. Our system covers
        resumption of interrupted work, file organisation, cited synthesis
        across documents, correspondence and voice, with code assistance as one
        capability among several and restricted to proposing a difference
        rather than writing to a file.
\end{itemize}

\section{What We Do Not Claim}

We do not solve approval fatigue. The SoK documents it as an open problem in
the exact mechanism our design depends upon, and our risk tiers reduce prompt
frequency without eliminating the effect. We have written it into our
evaluation as something to measure across repeated sessions rather than assert
away. We likewise inherit the limitations that a language model can produce
unfaithful summaries and unreliable citations; our bounded document set narrows
this exposure but does not remove it, and both are reported as known
limitations rather than as solved problems.

\section{Summary}

HULA shows the approval loop works but leaves the agent unbounded and
irreversible. The SoK shows confinement is the dependable defence but writes for
organisations with security staff. Our project sits in the gap: the validated
approval loop, inside a boundary the user grants, with reversal, for one
non-expert user on a personal machine.

\end{document}
